%! Right Angles Revisited --- Setting Up Orthogonality — Unit 4, Lesson 4.0 — Student Packet Cover Sheet
\documentclass[10pt]{article}
\usepackage{linalg-article}
\usepackage{linalg-boxes}
\usepackage{ltablex}
\keepXColumns

\begin{document}

% Full-bleed burgundy banner
\begin{tikzpicture}[remember picture, overlay]
  \fill[burgundy] (current page.north west)
    rectangle ([yshift=-0.9in]current page.north east);
\end{tikzpicture}
\vspace{-0.8in}

\begin{center}
  {\color{white}\textbf{\LARGE Linear Algebra}}\\[4pt]
  {\color{blushmid}\large\textbf{Unit 4 \quad Orthogonality}}\\[6pt]
  {\color{white}\textbf{\large Lesson 4.0 \quad Right Angles Revisited --- Setting Up Orthogonality}}
\end{center}

\vspace{0.22in}
\namedateperiod
\vspace{0.18in}

\begin{learningtargetbox}
\small
\begin{enumerate}[leftmargin=1.5em, itemsep=5pt]
  \item \ldots{} compute a \textbf{dot product}, find a vector's \textbf{length} $\|\mathbf v\|=\sqrt{\mathbf v\cdot\mathbf v}$, and shrink it to a \textbf{unit vector}.
  \item \ldots{} test whether two vectors are \textbf{perpendicular} using $\mathbf v\cdot\mathbf w=0$, and solve for an entry that makes them perpendicular.
  \item \ldots{} explain why the four subspaces come in \textbf{perpendicular pairs} and what Unit~4 is about (projection, least squares).
\end{enumerate}
\end{learningtargetbox}

\vspace{0.14in}

\begin{tocbox}
\small
\renewcommand{\arraystretch}{1.5}
\begin{tabularx}{\linewidth}{@{} c l X r @{}}
\rowcolor{blush}
\textbf{\#} & \textbf{Component} & \textbf{Description} & \textbf{Score} \\
\midrule
1 & Warm-Up        & Spiral review: dot products, length \& unit vectors, perpendicular check & \blank{1.2cm} \\
2 & Guided Notes   & The dot product, ``perpendicular $\Leftrightarrow$ dot $=0$,'' and the road ahead & \blank{1.2cm} \\
3 & Group Activity & Right Angles, Everywhere --- testing and building orthogonal vectors     & \blank{1.2cm} \\
4 & Exit Ticket    & Quick check on dot products, unit vectors, and orthogonality             & \blank{1.2cm} \\
5 & Homework       & Dot products, unit vectors, orthogonality, and a projection preview       & \blank{1.2cm} \\
\midrule
  & \multicolumn{2}{r}{\textbf{Total}} & \blank{1.2cm} \\
\end{tabularx}
\end{tocbox}

\vspace{0.10in}

\begin{remindbox}
\small The whole of Unit~4 rests on one fact from Lesson~1.2: two vectors are
\textbf{perpendicular} exactly when their \textbf{dot product is zero},
$\mathbf v\cdot\mathbf w=0$. A right angle becomes an equation. When a target is out of reach, the
\emph{closest} point is found by dropping a \textbf{perpendicular} --- that is
\emph{projection} (4.2) and \emph{least squares} (4.3), and it is where this unit is headed.
\end{remindbox}

\end{document}
